% !TEX root = chapter-standalone.tex

\section{Semigroup actions}
\label{sec:semigroup-actions}

A semigroup is like a monoid, but without a neutral element. The definition of a monoid action has a straightforward analogue for semigroups. 



\begin{ctdefinition}[Semigroup action]
    \label{def:semigroup-cov-action}
    An \maindef{action} of a \SY{semigroup} $\sgrpA = \tup{\sgrpAset, \mtimesof\sgrpA}$ on a set $\stateset$ is:

    \constit
    \begin{enumerate}
        \item A \SY{semigroup} morphism
              \begin{equation}
                  \label{eq:act-semigroup}
                  \act \colon \sgrpA \fto \Endof\stateset.
              \end{equation}
    \end{enumerate}
\end{ctdefinition}

\begin{ctdefinition}[Semigroup action, uncurried]
    \label{def:semigroup-uncurried-action}
    An \emph{uncurried action} of a \SY{semigroup} $\sgrpA = \tup{\sgrpAset, \mtimesof\sgrpA}$ on a set $\stateset$ is:

    \constit
    \begin{enumerate}
        \item a function
              \begin{equation}
                  \label{eq:semigroup-uncurried-action}
                  \act \colon \makecartprod{\sgrpAset, \stateset} \sto \stateset.
              \end{equation}
    \end{enumerate}
    \condit
    \begin{enumerate}
        \item for all $\statesetel \setin \stateset$ and $\sgrpelA, \sgrpelB \setin \sgrpAset$,
              \begin{equation}
                  \label{eq:semigroup-uncurried-action-1cond}
                  \act(\sgrpelA \mtimesof\sgrpA \sgrpelB, \statesetel) = \act(\sgrpelB, \act(\sgrpelA, \statesetel)).
              \end{equation}
    \end{enumerate}
\end{ctdefinition}

% \begin{ctdefinition}
%     [Semigroup contravariant action (preliminary version)]\label{def:semigroup-contra-action-prelim}
%     A \emph{semigroup \textbf{contravariant} action} of a \SY{semigroup}~$\sgrpA$ on a set~$\stateset$ is a map
%     \begin{equation}
%         \label{eq:ract}
%         \ract \colon \makecartprod{\sgrpAset, \stateset} \to \stateset
%     \end{equation}
%     such that, for all~$\statesetel \setin \stateset$,
%     \begin{equation}
%         \label{eq:ract1cond}
%         \ract(\monelb, \ract(\monela, \statesetel)) = \ract(\monelb \mthenof\sgrpA  \monela, \statesetel).
%     \end{equation}
% \end{ctdefinition}


\subsection{Actions of sets}

Following the pattern so far from monoids and semigroups, if we wanted to define an ``action of a set $\setA$ on a set $\stateset$'', then arguably such a thing would be defined simply as a function
\begin{equation}
\act \colon \setA \sto \Endof\stateset,
\end{equation}
without any further conditions. Although the term ``action'' is not usually used in this case, this situation is useful to also consider. 

If we apply (right) uncurrying to such a function $\act$, we obtain a function of the type
\begin{equation}
 \makecartprod{\setA, \stateset} \sto \stateset.
\end{equation}
We call this a \emph{parametrized endomorphism} (with parameter set $\setA$). The underlying dynamics of various flavors of dynamical systems are specified by such a function (see, for instance, \cref{exa:plant-trafo-semigroup}, \cref{def:LTI_syst_first_time}, and \cref{eq:moore-1}).

Given a function  
\begin{equation}
 \mora \colon \makecartprod{\setA, \stateset} \sto \stateset,
\end{equation}
we note that we can always build a monoid action of $\listsof{\setA}$ from it. Namely, the set $\listsof{\setA}$ of finite lists of elements of $\setA$ can be given a canonical monoid structure, with concatenation as composition and the empty list as neutral element (see \cref{exa:string-monoid}), and we obtain an action induced by $\mora$,
\begin{equation}
\act \colon \listsof{\setA} \fto \Endof\stateset,
\end{equation}
via
\begin{equation}
\act(\makelist{\setAel_1, \setAel_2, \dots, \setAel_k}) \definedas \mora(\setAel_1, -) \mthen \mora(\setAel_2, -) \mthen \cdots \mthen \mora(\setAel_k, -),
\qquad
\act(\emptylist) \definedas \catidat\stateset.
\end{equation}
In words: the action of a list processes its elements one after the other, from left to right; in uncurried form,
\begin{equation}
\act(\makelist{\setAel_1, \setAel_2, \dots, \setAel_k}, \statesetel) = \mora\bigl(\setAel_k, \dots \mora(\setAel_2, \mora(\setAel_1, \statesetel)) \dots\bigr).
\end{equation}
This is indeed a monoid homomorphism: the action of the concatenation of two lists is the action of the first list followed by the action of the second list, and the empty list acts as the identity.










